\documentclass[11pt]{article}
\usepackage[margin=0.85in]{geometry}\usepackage{amsmath,amssymb,booktabs,enumitem,xcolor,listings,fancyhdr,hyperref}
\definecolor{navy}{RGB}{24,55,95}\definecolor{green}{RGB}{25,100,70}\definecolor{codebg}{RGB}{246,248,250}\definecolor{codecomment}{RGB}{90,110,95}
\hypersetup{colorlinks=true,linkcolor=navy,urlcolor=navy}\pagestyle{fancy}\fancyhf{}\lhead{STAT452 Chapter 4}\rhead{R Exercises and Notes}\cfoot{\thepage}\setlength{\headheight}{14pt}
\lstset{language=R,basicstyle=\ttfamily\small,keywordstyle=\color{navy}\bfseries,commentstyle=\color{codecomment}\itshape,backgroundcolor=\color{codebg},frame=single,breaklines=true,columns=fullflexible,keepspaces=true,showstringspaces=false}
\newenvironment{exercise}[1]{\par\medskip\noindent\color{navy}\textbf{Exercise #1.}\color{black}\ }{\par\medskip}\newenvironment{answer}[1]{\par\medskip\noindent\color{green}\textbf{Solution #1.}\color{black}\ }{\hfill$\square$\par\medskip}
\title{Chapter 4: Logistic Regression and Classification in R\\\large Syntax Notes, Exercises, and Worked Solutions}\author{STAT452 -- Applied Statistics for Engineers and Scientists II}\date{}
\begin{document}\maketitle\tableofcontents\bigskip
\section{What this worksheet covers}
For a binary response $Y\in\{0,1\}$, linear regression can produce values outside $[0,1]$. Logistic regression models the log-odds: $\log\{\pi(x)/(1-\pi(x))\}=\beta_0+\beta_1x_1+\cdots+\beta_px_p$, where $\pi(x)=P(Y=1\mid x)$. R fits the coefficients by maximum likelihood with \texttt{glm(..., family=binomial)}. Exponentiating a coefficient gives an odds ratio.
\section{R syntax}
\begin{lstlisting}
fit <- glm(transmission ~ wt + hp, data=cars, family=binomial)
prob <- predict(fit, type="response")
pred <- ifelse(prob >= 0.5, "manual", "automatic")
exp(coef(fit))
anova(null_fit, fit, test="Chisq")
\end{lstlisting}
The 0.5 threshold is a decision rule, not a universal law. Report a confusion matrix and metrics such as accuracy, sensitivity, and specificity, preferably on held-out data.
\section{Exercises}
\begin{exercise}{1: Compare with linear regression} Regress binary \texttt{am} on \texttt{wt} and \texttt{hp}. Check whether fitted values lie in $[0,1]$.\end{exercise}
\begin{exercise}{2: Fit and interpret the logit model} Fit logistic regression. Compute odds ratios and explain coefficient signs while holding the other predictor fixed.\end{exercise}
\begin{exercise}{3: Classify observations} Convert probabilities to classes at threshold 0.5. Produce a confusion matrix and calculate accuracy.\end{exercise}
\begin{exercise}{4: Threshold trade-offs} Compare thresholds 0.3 and 0.5. Calculate sensitivity and specificity.\end{exercise}
\begin{exercise}{5: Validate predictions} Split into training and test observations. Fit on training data and evaluate on test data. Explain optimistic training accuracy.\end{exercise}
\begin{exercise}{6: Compare models} Use a likelihood-ratio test and AIC to compare intercept-only and two-predictor models.\end{exercise}

\section{Worked solutions}
\begin{answer}{1} Linear regression is not guaranteed to respect probability bounds; 
    \newline
    inspect \texttt{range(predict(linear\_binary))}. This motivates the logistic link.\end{answer}
\begin{answer}{2} Use \texttt{exp(coef(fit))}. An odds ratio above one raises the odds of the modeled success class; below one lowers the odds, conditional on other predictors.\end{answer}
\begin{answer}{3} A confusion matrix separates correct and incorrect classifications. Accuracy is $(TP+TN)/n$, but it can hide unequal costs or class imbalance.\end{answer}
\begin{answer}{4} Lowering the threshold generally increases sensitivity and decreases specificity. Choose it based on the cost of false negatives versus false positives.\end{answer}
\begin{answer}{5} Test observations were not used during estimation, so test performance better reflects generalization. With 32 cars, treat this split as practice rather than a definitive performance claim.\end{answer}
\begin{answer}{6} The likelihood-ratio test compares nested models using change in deviance. AIC balances likelihood and parameter count; lower AIC is preferred among models compared.\end{answer}
\section*{Checklist}\texttt{glm -> predict(type=response) -> threshold -> confusion matrix -> held-out validation}.\end{document}

